\noindent\textbf{Elección del procedimiento.} En el marco aproximadamente inercial de la pista, el ave pasa de velocidad despreciable a $V$. El impulso cumple $\bar F\Delta t=\Delta(mv)=mV$. La fuerza que el ave ejerce sobre el cristal es opuesta a la que este ejerce sobre el ave.

\begin{enumerate}

\item Calcular la variación de momento, no la energía cinética.

\[\fitmath{\mathrm{dpave}=m\cdot V}\]

\[\fitmath{\mathrm{dpave}=(1.8)\cdot (80)\approx 144\,\mathrm{kg\,m/s}}\]

\item Convertir 2 ms a 0.002 s y dividir el impulso por la duración.

\[\fitmath{\mathrm{Fmean}=\frac{\mathrm{dpave}}{\mathrm{dt}}}\]

\[\fitmath{\mathrm{Fmean}=\frac{(144)}{(0.002)}\approx 72000\,\mathrm{N}}\]

\end{enumerate}

\noindent\textbf{Contraste y comparación.} Esta información determina la fuerza media. No determina la fuerza máxima ni cuántas veces supera a la media: haría falta la historia temporal del impacto.
